\documentclass[11pt]{article}
\usepackage[T1]{fontenc}
\usepackage{amsmath,amssymb,amsthm}
\usepackage[margin=1in]{geometry}
\usepackage{booktabs}
\usepackage{microtype}
\usepackage[hidelinks,pdfusetitle]{hyperref}

\newtheorem{theorem}{Theorem}
\newtheorem{lemma}[theorem]{Lemma}
\newcommand{\CC}{\mathbb C}
\newcommand{\NN}{\mathbb N}

\title{The Borel-fixed colength-15 tangent maximum on \\
  $\operatorname{Hilb}^{15}(\mathbb A^3)$}
\author{Xiaojian Zeng\\
PBC School of Finance, Tsinghua University, 100190\\
\texttt{xiaojianzeng@foxmail.com}}
\date{}

\begin{document}
\maketitle

\begin{abstract}
For the fixed Borel order $x<y<z$, we classify the zero-dimensional
Borel-fixed monomial ideals of colength $15$ in $\CC[x,y,z]$ that maximize
$\dim_{\CC}\operatorname{Hom}(I,R/I)$.  An exhaustive exact census contains
$107$ ideals.  The maximum tangent dimension is $81$, attained by exactly
four ideals, all having smallest pure-power exponent $m_1=3$.  The largest
dimension when $m_1\leq2$ is $69$ and has a unique equality case.  This
settles the colength-15 Borel-fixed cell of the necessary-condition
conjecture, but it makes no assertion about the maximum over non-Borel
points of the Hilbert scheme.
\end{abstract}

\section{Statement and context}

Put $R=\CC[x,y,z]$, with the variable order $x<y<z$.  For a
zero-dimensional Borel-fixed monomial ideal $I$, write
\[
 T(I)=\dim_{\CC}\operatorname{Hom}_R(I,R/I)
\]
and let $m_i(I)$ be the exponent of the least pure power in the $i$th
variable that belongs to $I$.  Thus $m_1\leq m_2\leq m_3$.

Rezaee formulated necessary and sufficient conjectural criteria for maximum
tangent dimension and gave explicit candidate families
\cite{Rezaee2025}.  At colength $15$, her family includes an ideal with
tangent dimension predicted to be $81$.  Ascott, Rezaee, and Zhou later
reported the exact stratum maxima $45,69,81$ for $m_1=1,2,3$ respectively
\cite{AscottRezaeeZhou2026}.  Their table establishes the numerical
comparison relevant to the necessary condition, but does not print a
complete colength-15 atlas or the list of all equality cases.

Rezaee and Saberi prove the necessary condition for three explicit families
of colengths \cite{RezaeeSaberi2026}.  In the first interval
$10\leq l<20$, substituting $k=3$ in those families covers
$10,11,12,13,14,16,17,18,19$; hence $15$ is the one colength in that
interval outside their theorem.  The contribution here is an exact finite
resolution of that Borel-fixed cell and a classification of every equality
case.  No claim of priority beyond the cited, bounded comparison is made.

\begin{theorem}\label{thm:main}
Among all zero-dimensional Borel-fixed monomial ideals of colength $15$ in
$R$, one has
\[
 \max_I T(I)=81.
\]
Exactly four ideals attain the maximum:
\begin{align*}
I_1={}&(y^3,xyz,xy^2,x^2z,x^2y,x^3,yz^3,y^2z^2,xz^3,z^5),\\
I_2={}&(y^2z,y^3,xyz,xy^2,x^2z,x^2y,x^3,yz^3,xz^3,z^6),\\
I_3={}&(xyz,xy^2,x^2z,x^2y,x^3,z^4,yz^3,y^2z^2,y^3z,y^4,xz^3),\\
I_4={}&(xz^2,xyz,xy^2,x^2z,x^2y,x^3,yz^3,y^2z^2,y^3z,y^4,z^5).
\end{align*}
Their pure-power tuples are, respectively,
\[
 (3,3,5),\qquad(3,3,6),\qquad(3,4,4),\qquad(3,4,5).
\]
In particular every Borel-fixed maximizer has $m_1=3$.  Among the ideals
with $m_1\leq2$, the maximum is $69$ and is attained uniquely by
\[
 (x^2,xyz,xy^2,yz^3,y^2z^2,y^3z,y^4,xz^3,z^5),
\]
whose pure-power tuple is $(2,4,5)$.
\end{theorem}

\section{Complete finite encoding}

For a monomial ideal $I$, let
\[
 O(I)=\{(a,b,c)\in\NN^3:x^ay^bz^c\notin I\}.
\]
Colength $15$ says that $|O(I)|=15$.  The set $O(I)$ is coordinatewise
downward closed.  With the chosen Borel order, strong stability of $I$ is
equivalent to the following additional closures of its standard set:
whenever $a_i>0$ and $i<j$, the vector obtained by moving one unit from
coordinate $i$ to coordinate $j$ also belongs to $O(I)$.

The census was constructed in two exact ways.  The first begins with the
origin, successively adds every available outer corner, and canonicalizes
the resulting coordinate downsets; at size $15$ it retains precisely those
sets satisfying all three Borel moves.  This is complete because deleting a
coordinatewise maximal point from a finite downset gives the reverse
induction, while the final closure test is equivalent to strong stability.

The second construction assigns to each exponent vector its least closure
under coordinate predecessors and Borel right shifts, then recursively forms
unions of these closures.  Every closed set is the union of the one-point
closures of its members.  A point of total degree $d$ forces a coordinate
chain of $d+1$ points, so degree at most $14$ is sufficient for a set of
size $15$.  Pruning unions only after their size exceeds $15$ therefore
cannot discard a target.  The two constructions give the same $107$
canonical exponent sets, not merely the same cardinality.  Minimal
generators are recovered as the minimal boundary of each complement, and
the pure powers are read from the three coordinate axes.

\section{Exact tangent calculation}

Fix one census ideal, put $A=R/I$, and let $g_1,\ldots,g_s$ be its minimal
monomial generators.  A homomorphism $I\to A$ is specified by images
$a_i\in A$ subject to the relations among the $g_i$.  The pairwise Taylor
relations \cite{Taylor1966}
\[
 \frac{\operatorname{lcm}(g_i,g_j)}{g_i}e_i-
 \frac{\operatorname{lcm}(g_i,g_j)}{g_j}e_j,
 \qquad i<j,
\]
generate the kernel of $\bigoplus_i Re_i\to I$.  Expanding every $a_i$ in
the $15$ standard monomials therefore produces an integer matrix whose
nullity is $T(I)$.  Exact rational row reduction is sufficient because an
integer matrix has the same rank over $\mathbb Q$ and $\CC$.

As a structurally separate check, each ideal was also given its
Eliahou--Kervaire stable presentation \cite{EliahouKervaire1990}.  After
expansion in the standard basis, every surviving coefficient equation in
this presentation either pins one coefficient to zero or equates two
coefficients.  Thus its nullity is the number of connected equality-graph
components not pinned to zero.  This component calculation agrees with the
Taylor-matrix rank for every one of the $107$ ideals.  A separately written
verifier reconstructs the stable factorizations and repeats this calculation
for the four maximizers and the closest $m_1\leq2$ case.  The same general
certificate pattern---canonical exhaustive input, exact arithmetic, and an
independent reconstructor---has been used in another finite monomial-algebra
problem \cite{Zeng2026SSRN7385138}; here both presentations additionally agree on
the full atlas.

The resulting distribution is as follows.
\begin{center}
\small
\begin{tabular}{@{}rrrrrrrrrrrrrr@{}}
\toprule
$T$&45&51&53&55&57&59&61&63&65&67&69&75&81\\
\midrule
count&27&8&11&5&8&10&7&10&4&5&1&7&4\\
\bottomrule
\end{tabular}
\end{center}
The counts sum to $107$.  Extracting the four records in the last column
gives exactly $I_1,\ldots,I_4$.  Filtering the same complete table by
$m_1\leq2$ leaves $96$ records; its last nonempty column is $69$ and contains
only the ideal stated in Theorem~\ref{thm:main}.  This proves the theorem.

\section{Scope and reproducibility}

The theorem is only a classification within the Borel-fixed points.  It
does not bound tangent dimensions at non-Borel points and therefore does not
classify the most singular points of the full
$\operatorname{Hilb}^{15}(\mathbb A^3)$.

The accompanying archive contains the canonical $107$-record atlas, the
complete tangent table, the two census implementations, the Taylor and
Eliahou--Kervaire computations, and the standalone extrema verifier.  All
enumeration predicates and rank decisions use integer or rational
arithmetic; no floating-point rank or probabilistic prime is used.

The author used OpenAI Codex to assist with literature discovery, program
development, proof organization, exact-output auditing, and manuscript
drafting.  The mathematical claim rests on the stated finite arguments and
replayable exact certificates, not on model output.  The software routines
were checked through the two complete census constructions, two syzygy
presentations on all records, and the separately implemented extrema
reconstruction.

\bibliographystyle{plain}
\bibliography{references}

\end{document}
